% Generated by results/thesis_tables.py - do not edit by hand.
\begin{table}[H]
\centering
\footnotesize
\begin{tabular}{|>{\raggedright\arraybackslash}p{4.2cm}|r|r|>{\centering\arraybackslash}p{1.45cm}|>{\centering\arraybackslash}p{1.45cm}|>{\centering\arraybackslash}p{1.45cm}|>{\centering\arraybackslash}p{1.45cm}|}
\hline
Σύνολο & Πακέτα & Χωρίς chains & Πλήρες σύστημα & Μόνο δείγματα & Μόνο πολιτικές & Χωρίς RAG \\
\hline
Χωρίς συγκάλυψη (κακόβουλα) & 97 & 0 & 94/97 & 92/97 & 93/97 & 92/97 \\
Επίπεδο A (κακόβουλα) & 97 & 12 & 83/85 & 83/85 & 82/85 & 84/85 \\
\quad χωρίς αρχεία που δεν αναλύονται & 70 & 1 & 68/69 & 68/69 & 67/69 & 69/69 \\
Επίπεδο B (κακόβουλα) & 97 & 0 & 97/97 & 95/97 & 95/97 & 97/97 \\
\hline
Χωρίς συγκάλυψη (νόμιμα) & 99 & 0 & 3/99 & 3/99 & 3/99 & 1/99 \\
Επίπεδο A (νόμιμα) & 99 & 21 & 30/78 & 20/78 & 20/78 & 14/78 \\
\quad χωρίς αρχεία που δεν αναλύονται & 6 & 0 & 0/6 & 0/6 & 0/6 & 0/6 \\
Επίπεδο B (νόμιμα) & 99 & 0 & 96/99 & 96/99 & 90/99 & 88/99 \\
\hline
\end{tabular}
\caption{Τα αντίγραφα με συγκάλυψη (Intensio-Obfuscator: επίπεδο A μετονομασία και νεκρός κώδικας, επίπεδο B επιπλέον κωδικοποίηση κάθε αρχείου σε hex που εκτελείται με \texttt{exec}) δίπλα στα ίδια πακέτα χωρίς συγκάλυψη. Σε κάθε κελί: πακέτα που χαρακτηρίστηκαν \texttt{MALICIOUS} / πακέτα με Behavior Chains.}
\label{tab:obfuscation}
\end{table}
